\section{定積分の基本性質}\label{節:定積分性質}
以下では，特に断らない限り，被積分関数は扱う閉区間で連続であるとする。まず，積分の端点が同じ場合と，順序が逆の場合を
\[
 \int_a^a f(x)\dd x=0,\qquad
 \int_b^a f(x)\dd x=-\int_a^b f(x)\dd x\quad(a<b)
\]
と定める。積分の向きが反対になると符号が反対になる。この約束により，区間をつなぐ式を端点の大小にかかわらず統一して書ける。

\begin{theorem}{線形性と区間の加法性}\label{定理:積分基本性質}
$\alpha,\beta$ を実数とする。関係する区間で $f,g$ が連続なら，
\begin{align*}
 \int_a^b\{\alpha f(x)+\beta g(x)\}\dd x
 &=\alpha\int_a^b f(x)\dd x+\beta\int_a^b g(x)\dd x,\\
 \int_a^b f(x)\dd x
 &=\int_a^c f(x)\dd x+\int_c^b f(x)\dd x
\end{align*}
が成り立つ。後者は $a,c,b$ の順序によらない。
\end{theorem}

線形性は，同じ分割と代表点を使った有限和について
\[
 \sum_k\{\alpha f(\xi_k)+\beta g(\xi_k)\}\Delta x_k
 =\alpha\sum_k f(\xi_k)\Delta x_k
  +\beta\sum_k g(\xi_k)\Delta x_k
\]
と分け，極限をとれば得られる。区間の加法性は，$a<c<b$ のとき，分割に $c$ を含めて和を二つに分ければ分かる。端点の順序が異なる場合は，逆向きの積分の定義から同じ式になる。積分区間を分けても，被積分関数を勝手に変えているわけではない。

\begin{theorem}{大小関係と絶対値の評価}\label{定理:積分評価}
$a\le b$ とする。区間 $[a,b]$ で連続な関数について，次が成り立つ。
\begin{align*}
 f(x)\le g(x)\ (a\le x\le b)
 &\quad\Longrightarrow\quad
 \int_a^b f(x)\dd x\le\int_a^b g(x)\dd x,\\
 \left|\int_a^b f(x)\dd x\right|
 &\le\int_a^b|f(x)|\dd x.
\end{align*}
特に，$m\le f(x)\le M$ なら，
\[
 m(b-a)\le\int_a^b f(x)\dd x\le M(b-a).
\]
\end{theorem}

最初の性質は，各小区間の幅が正であり，リーマン和の各項について大小関係が保たれることから分かる。二つ目は $-|f|\le f\le|f|$ を積分すれば得られる。最後の評価では，定数関数の積分が高さと区間の長さの積になることを使った。これらの不等式は，端点を $a\le b$ の順に置いていることが前提である。

\begin{example}{原始関数を探さずに範囲を決める}
$\int_0^1\frac1{1+x^2}\dd x$ が $\frac12$ 以上 $1$ 以下であることを示せ。
\end{example}
$0\le x\le1$ では $1\le1+x^2\le2$ であり，すべて正なので逆数の大小は逆になる。
\[
 \frac12\le\frac1{1+x^2}\le1.
\]
区間の長さは $1$ だから，これを積分して
\[
 \frac12\le\int_0^1\frac1{1+x^2}\dd x\le1
\]
を得る。積分の値を正確に計算できなくても，関数の高さを比較すれば範囲を保証できる。この例の原始関数は\chapterref{章:置換の背景}で扱う。

\subsection{関数の平均値}
$a<b$ のとき，
\[
 f_{\mathrm{av}}=\frac1{b-a}\int_a^b f(x)\dd x
\]
を $f$ の区間 $[a,b]$ における\textbf{平均値}という。等分割のリーマン和を区間の長さで割ると $\frac1n\sum_{k=1}^{n}f(\xi_k)$ になるので，有限個の値の平均を細かくしていったものと解釈できる。微分で用いた平均変化率とは別の量である。

$f$ の最小値を $m$，最大値を $M$ とすると，積分の評価から $m\le f_{\mathrm{av}}\le M$ である。連続関数は最小値と最大値の間の値をすべてとるから，ある $c\in[a,b]$ で
\[
 \int_a^b f(x)\dd x=f(c)(b-a)
\]
が成り立つ。これを\textbf{積分の平均値の定理}という。高さを一定の $f(c)$ にした長方形と，元の関数の定積分が等しくなるという意味である。

%
